\section{Conclusion and Future Work}
\label{sec:conclusion}

In this paper, we presented \textbf{ChaosMerge} (Chaos-Theoretic Attractor Merging), a novel and highly effective framework for dynamic parameter-space model merging. By rejecting the traditional view that network depth is simply a feed-forward sequence of flat computations, we conceptualize layer depth as the temporal progression of a non-linear chaotic dynamical system. We formulate the weight-merging coefficients as the trajectory of a Coupled Map Lattice (CML) driven by a chaotic Logistic Map, steered by input-specific unit-sphere phase projections and logit perturbations.

To resolve the fundamental optimization challenges of deep recurrent chaotic systems, we introduced the \textbf{Gated Coupled Map Lattice (G-CML)}, which incorporates learned layer-wise gating as residual skip-connections. G-CML successfully tames catastrophic gradient explosion (limiting the exponential $4^{14}$ multiplier), enabling robust optimization via standard first-order gradient descent. Furthermore, we addressed the batch-averaging contradiction by employing \textbf{Task-Specific Dynamic Routing}, which extracts task-specific merging coefficients directly and prevents individual task signatures from being washed out in heterogeneous batches.

Empirical evaluations across visual classification tasks using a ViT backbone demonstrate that ChaosMerge achieves highly competitive results. With only \textbf{384 parameters}, ChaosMerge significantly outperforms competitive static baselines (such as Task Arithmetic) and unsupervised test-time static merging (such as AdaMerging), while achieving performance competitive with supervised static tuning (OFS-Tune) and unconstrained, over-parameterized dynamic routers with nearly $30\times$ more parameters, validating the exceptional regularizing and separating power of stabilized chaotic dynamics in parameter space.

\subsection{Future Directions}
Our findings lay out a clear roadmap for future research in non-linear parameter fusion:
\begin{itemize}
    \item \textbf{Outstanding Scale-Invariant Complexity:} A major highlight of G-CML is its outstanding parameter scalability. Since G-CML's parameter complexity $\mathcal{O}(LK)$ is determined solely by the number of experts $K$ and layer groups $L$, and is completely decoupled from the backbone's internal hidden dimension (via spherical random projection), the routing footprint remains extremely small. For example, scaling G-CML to a massive 32-layer LLM (like Llama-3-8B) with $K=8$ experts would require fewer than \textbf{2,000 trainable parameters} total—a negligible footprint compared to standard adapters or full parameter tuning, providing a powerful, parameter-efficient case for modern large-scale AI.
    \item \textbf{Resolving the Unsupervised Clustering Bottleneck:} While our empirical validation highlighted the fragility of on-the-fly unsupervised $K$-means clustering in the projected phase space (low purity and catastrophic error propagation), future work will investigate robust cluster-number estimation (e.g., Dirichlet Process mixture models) and lightweight multi-centroid routing to support fully task-agnostic, mixed-task deployments on edge devices.
    \item \textbf{Alternative Dynamical Maps and Continuous-Time Flows:} We plan to explore other non-linear maps, such as the H{\'e}non map or physical systems (e.g., thermodynamic spin glasses), which may offer distinct geometric properties while retaining rich self-organizing representation spaces. Furthermore, we aim to extend our temporal view of depth to continuous-time neural ordinary differential equations (Neural ODEs), modeling parameter trajectories as continuous flows on Riemannian manifolds.
\end{itemize}
